
Large Language Models (LLMs) such as GPT \cite{openai_gpt4_2023} and LLaMA \cite{touvron_llama_2023} are now deployed worldwide and used in many languages. As these systems mediate information access and public discourse, a central concern is whether they present the same opinions and reasoning to users across countries and languages—or whether they exhibit ideological or cultural biases \cite{weidinger_risks_2021,pmlr-v202-santurkar23a}. Understanding these effects is essential for trustworthy, fair, and transparent AI \cite{navigating-llm-ethics}.

Existing audits have shown that LLMs can encode biases along several axes, including social, linguistic, and representational dimensions \cite{kurita2019bias,stereoset2021,detecting2023}. However, most empirical studies remain limited to a single country, a single demographic axis, or rely heavily on manual prompt probes and human judgments \cite{gender_gpt2023}. Less attention has been paid to cross-country and cross-language ideological behavior in contemporary chat models, and to disentangling the influence of role framing (e.g., “answer as a citizen of X”) from language framing (answering in X’s native language). Recent work has begun to examine political and national biases in LLMs \cite{silence_llms2023,nationality_bias2013,russia_china_war2024}, but reproducibility remains a challenge, since manual inspection is often subjective and difficult to replicate. This motivates the need for a multilingual and reproducible evaluation that targets both stance and reasoning on policy-relevant topics, while reducing dependence on ad-hoc human judging.

The objective of this thesis is to build such a framework. My approach is to evaluate ideological bias across roles and languages in a way that combines stance alignment with semantic similarity of justifications \cite{detecting2023}. In addition, I compare whether different model families—GPT and LLaMA—exhibit similar or diverging bias patterns, thereby situating model behavior within a broader comparative perspective \cite{silence_llms2023}. This dual focus on stance and reasoning aims to provide a more fine-grained picture than agreement counts alone, while ensuring reproducibility.

Taken together, these motivations lead to four guiding research questions:
\begin{enumerate}
    \item Do LLMs exhibit ideological bias across countries and policy topics?
    \item Are model outputs more sensitive to role framing or to language framing?
    \item How similar are role-based and native-language responses for the same country?
    \item Are the observed effects consistent across model families (GPT vs.\ LLaMA)?
\end{enumerate}

\section{Thesis Structure}

This thesis is organized into six chapters, followed by the references and an appendix. Chapter~2 surveys the state of the art, introducing core ideas in Natural Language Processing from static to contextual embeddings, explaining how these advances led to Large Language Models, reviewing bias in NLP and LLMs, and concluding with a gap analysis. Chapter~3 presents the methodology: the dataset and preprocessing (including topic modeling), the prompting and response-collection setup for role-based and language-based conditions, and the analysis metrics—stance alignment via agreement ratios and embedding-based comparisons of justifications. Chapter~4 reports the results, including topic-wise agreement ratios for GPT-4o-mini and LLaMA~3.1-70B, embedding-similarity analyses summarized with bar charts and heatmaps, and an error/neutrality analysis. Chapter~5 distills the conclusions, relating the findings to the research questions and prior work. Chapter~6 discusses limitations and outlines directions for future research. The \textit{References} list all cited sources, and the \textit{Appendix} provides supporting material such as translation maps, t-SNE visualizations of topic embeddings, and aggregated radar plots for cross-condition comparisons.
